\section{Model cards and equations}
\label{sec:s-models-v2}

\subsection{Node classifiers}

The feature-only MLP does not consume adjacency and therefore serves as a negative control for the graph-only visibility intervention. The GCN update is
\begin{equation}
 H^{(\ell+1)}
 =\sigma\!\left(\widehat D^{-1/2}\widehat A\widehat D^{-1/2}
 H^{(\ell)}W^{(\ell)}\right),
\end{equation}
with self-loops in \(\widehat A\)~\cite{kipf2017gcn}. Mean GraphSAGE uses
\begin{equation}
 h_v^{(\ell+1)}
 =\sigma\!\left(W^{(\ell)}
 [h_v^{(\ell)}\Vert |N(v)|^{-1}\!\sum_{u\in N(v)}h_u^{(\ell)}]\right)
\end{equation}
and therefore depends differently on available neighborhoods~\cite{hamilton2017graphsage}.

\begin{table}[tbp]
\centering
\footnotesize
\setlength{\tabcolsep}{3.2pt}
\renewcommand{\arraystretch}{1.12}
\caption{Node-classifier cards for Elliptic and DGraphFin.}
\label{tab:s-node-models}
\begin{tabularx}{\textwidth}{@{}>{\raggedright\arraybackslash}p{0.16\textwidth} Y Y >{\raggedright\arraybackslash}p{0.22\textwidth}@{}}
\toprule
\textbf{Model} & \textbf{Computation} & \textbf{Scientific role} & \textbf{Configuration} \\
\midrule
MLP & node features only & negative control for graph-visibility intervention & Elliptic h256/l3; DGraphFin h64/l2 \\
GCN & normalized neighborhood aggregation & tests spectral-style smoothing under visibility change & Elliptic h256/l3; DGraphFin h64/l2 \\
GraphSAGE & mean neighborhood aggregation & tests inductive aggregation under visibility change & Elliptic h256/l3; DGraphFin h64/l2 \\
\bottomrule
\end{tabularx}
\vspace{2pt}
\begin{minipage}{\textwidth}\footnotesize\emph{Scope and reading.} The MLP is the graph-visibility negative control. The grid does not establish class-wide behavior for modern temporal or heterogeneous GNNs.\end{minipage}
\end{table}


The card makes the architecture control explicit: Elliptic uses h256/l3 while DGraphFin uses h64/l2. The grid studies three mostly static families and does not support a class-wide conclusion about temporal-state, transformer, or heterogeneous GNNs.

\subsection{IBM transaction classifiers}

Logistic regression and histogram gradient boosting operate on transaction attributes without account-graph message passing~\cite{friedman2001gbm}. The SAGE-derived edge classifiers first form one-hop training-history summaries
\begin{equation}
 \bar x_{N(v)}
 =\frac{1}{\max(1,|N_{\mathrm{tr}}(v)|)}
 \sum_{u\in N_{\mathrm{tr}}(v)}x_u,
 \qquad
 \widetilde x_v=[x_v\Vert\bar x_{N(v)}],
\end{equation}
then score transaction \((u,v)\) with
\begin{equation}
 q_{uv}=\sigma\!\left(f_\theta(
 [\widetilde x_u\Vert\widetilde x_v\Vert a_{uv}])\right).
\end{equation}
They are static one-hop edge classifiers, not stacks of sampled GraphSAGE convolutions.

The feasible GINE candidate uses one edge-conditioned update,
\begin{equation}
 h_v'=\operatorname{MLP}\!\left((1+\epsilon)h_v+
 \sum_{u\in N(v)}\operatorname{ReLU}(h_u+\phi(a_{uv}))\right),
\end{equation}
followed by an edge head~\cite{hu2020pretraining}.

\begin{table}[tbp]
\centering
\footnotesize
\setlength{\tabcolsep}{3.2pt}
\renewcommand{\arraystretch}{1.12}
\caption{IBM transaction-classifier cards.}
\label{tab:s-ibm-models}
\begin{tabularx}{\textwidth}{@{}>{\raggedright\arraybackslash}p{0.20\textwidth} Y Y >{\raggedright\arraybackslash}p{0.20\textwidth}@{}}
\toprule
\textbf{Model} & \textbf{Input/computation} & \textbf{Scientific role} & \textbf{Configuration} \\
\midrule
Logistic regression & standardized transaction features & linear non-graph baseline & balanced classes; max\_iter=500 \\
Histogram gradient boosting & transaction features & nonlinear tabular baseline & 160 iterations; learning rate 0.06 \\
GraphSAGE-derived edge classifier h32 & training-history endpoint/neighborhood summaries plus edge features & small IBM graph baseline & 20 epochs; batch 65,536 \\
Edge-aware GraphSAGE-derived edge classifier h64 & training-history endpoint/neighborhood summaries plus edge features & reference for graph-feature/construction ablations & 30 epochs; batch 32,768 \\
GINE h64 & one GINE layer and edge head & edge-conditioned message-passing architecture & 30 epochs; batch 32,768; Small only \\
\bottomrule
\end{tabularx}
\vspace{2pt}
\begin{minipage}{\textwidth}\footnotesize\emph{Scope and reading.} The h32/h64 systems are static one-hop GraphSAGE-derived edge classifiers. Medium GINE is resource-blocked and has no performance value.\end{minipage}
\end{table}


GINE h64 is measured only on Small. The two Medium lanes exhausted a T4 and receive no performance score. The model card preserves that feasibility condition alongside the computational definition.
